% Part: lambda-calculus
% Chapter: syntax
% Section: substitution

\documentclass[../../../include/open-logic-section]{subfiles}

\begin{document}

\olfileid{lam}{syn}{sub}
\olsection{સ્થાનાપન્ન}

\begin{explain}
મુક્ત ચલો પરિવેશના ચલો તરફ નિર્દેશ કરે છે; તેથી મુક્ત ચલની
જગ્યાએ કોઈ નિશ્ચિત મૂલ્ય વાપરવું અર્થપૂર્ણ છે. દાખલા તરીકે,
$\lambd[x][f x]$માંના~$f$ની જગ્યાએ નિશ્ચિત પદ, જેમ કે
ઓળખ-વિધેય~$\lambd[y][y]$, મૂકવું હોય. પરિણામે
$\lambd[x][(\lambd[y][y]) x]$ મળે છે. મુક્ત ચલોની જગ્યાએ
લૅમ્બડા પદો મૂકવાની પ્રક્રિયા \emph{સ્થાનાપન્ન} કહેવાય છે.
\end{explain}

\begin{defn}[સ્થાનાપન્ન] \ollabel{defn:substitution}
  પદ~$M$માં ચલ~$x$ની જગ્યાએ પદ~$N$નું \emph{સ્થાનાપન્ન},
  $\Subst{M}{N}{x}$, આગમનથી આ રીતે વ્યાખ્યાયિત થાય છે:
  \begin{enumerate}
    \item $\Subst{x}{N}{x} = N$. \ollabel{defn:substitution-1}
    \item $\Subst{y}{N}{x}  = y$ જો $x \neq y$ હોય.
      \ollabel{defn:substitution-2}
    \item $\Subst{PQ}{N}{x}  = (\Subst{P}{N}{x}) (\Subst{Q}{N}{x})$.
      \ollabel{def:substitution-3}
    \item $\Subst{(\lambd[y][P])}{N}{x} = \lambd[y][\Subst{P}{N}{x}]$,
      જો $x \neq y$ અને $y \notin \FV{N}$ હોય તથા અંદરનું
      સ્થાનાપન્ન વ્યાખ્યાયિત હોય. જો $x=y$ હોય અથવા
      $x\notin\FV{P}$ હોય તો આખું અમૂર્તીકરણ યથાવત રહે છે;
      બાકીના કિસ્સામાં સ્થાનાપન્ન અવ્યાખ્યાયિત છે.
      \ollabel{defn:substitution-4}
  \end{enumerate}
\sourcecorrection{OLLAM-012}{સ્થિર અંગ્રેજી મૂળનો ચોથો નિયમ
બાકીના બધા કિસ્સાઓને અવ્યાખ્યાયિત કહે છે, જ્યારે તરત પછીનું
ઉદાહરણ બદ્ધ ચલને બદલી ન નાખતાં અમૂર્તીકરણ યથાવત રાખવાનું કહે
છે. અહીં બાંધનાર ચલ જ બદલવાનું હોય, અથવા બદલવાના ચલની
મુક્ત ઘટના અંદરના પદમાં ન હોય, તો અમૂર્તીકરણ યથાવત રાખ્યું
છે; આથી નિષ્પ્રયોજન સ્થાનાપન્નમાં પણ ખોટો બંધન-અવરોધ ઊભો
થતો નથી.}
\end{defn}

\begin{explain}
\olref{defn:substitution}\olref{defn:substitution-4}માં
$x \neq y$ની શરત એટલા માટે છે કે~$M$માં ચલ~$x$ની
\emph{બદ્ધ} ઘટનાઓને~$N$થી બદલવી નથી. દાખલા તરીકે,
$\Subst{\lambd[x][x]}{y}{x}$ ગણીએ તો પરિણામ
$\lambd[x][y]$ નહીં પરંતુ~$\lambd[x][x]$ જ હોવું જોઈએ.

$\lambd[y][P]$માં~$x$ની જગ્યાએ~$N$ મૂકતાં આપણે
$y \notin \FV{N}$ પણ માગીએ છીએ. દાખલા તરીકે,
$\lambd[y][x]$માં~$x$ની જગ્યાએ~$y$ મૂકી શકાતું નથી;
અર્થાત્~$\Subst{\lambd[y][x]}{y}{x}$ અવ્યાખ્યાયિત છે.
એમ કરવાથી~$\lambd[y][y]$ મળશે, જે આર્ગ્યુમેન્ટ સ્વીકારીને
તે જ પાછું આપતું વિધેય છે. પરંતુ~$\lambd[y][x]$ હંમેશા
પદ~$x$ (અથવા~$x$ જેનો નિર્દેશ કરે તે) પાછું આપતું વિધેય છે.
આથી આપણને ખરેખર એવું વિધેય જોઈએ જે આર્ગ્યુમેન્ટ સ્વીકારે,
તેને અવગણે અને પરિવેશના ચલ~$y$ને પાછું આપે. તે યોગ્ય રીતે
કરવા પહેલાં બદ્ધ ચલ~$y$નું ``નામ બદલવું'' પડે.
\end{explain}

\begin{prob}
  નીચેના સ્થાનાપન્નોનાં પરિણામો શું છે?
  \begin{enumerate}
  \item $\Subst{\lambd[y][x(\lambd[w][vwx])]}{(uv)}{x}$
  \item $\Subst{\lambd[y][x(\lambd[x][x])]}{(\lambd[y][xy])}{x}$
  \item $\Subst{y(\lambd[v][xv])}{(\lambd[y][vy])}{x}$
  \end{enumerate}
\end{prob}

\begin{thm} \ollabel{thm:notinfv}
  $x \notin \FV{M}$ હોય અને સ્થાનાપન્ન વ્યાખ્યાયિત હોય તો
  $\FV{\Subst{M}{N}{x}} = \FV{M}$.
\end{thm}

\begin{proof}
  $M$ના રચનાક્રમ પર આગમનથી.
  \begin{enumerate}
  \item $M$ ચલ હોય: અભ્યાસપ્રશ્ન.
  \item $M$ એ~$(PQ)$ સ્વરૂપનું હોય: અભ્યાસપ્રશ્ન.
  \item $M$ એ~$\lambd[y][P]$ સ્વરૂપનું હોય. જો બદલવાનો ચલ
    બાંધનાર ચલ જ હોય તો નિયમ અનુસાર પદ યથાવત રહે છે અને
    દાવો તરત મળે છે. અન્યથા,
    $\Subst{\lambd[y][P]}{N}{x}$ વ્યાખ્યાયિત હોવાથી તે
    $\lambd[y][\Subst{P}{N}{x}]$ છે. તેથી
    $\Subst{P}{N}{x}$ વ્યાખ્યાયિત છે; વળી $x \neq y$
    અને~$x \notin \FV{P}$. પછી:
    \begin{multline*}
      \FV{\Subst{\lambd[y][P]}{N}{x}} \\
    \begin{aligned}
      & = \FV{\lambd[y][\Subst{P}{N}{x}]} &&
      \text{\olref{defn:substitution-4} મુજબ}\\
      & = \FV{\Subst{P}{N}{x}} \setminus \{y\} &&
      \text{\olref[fv]{def:fv}\olref[fv]{def:fv2} મુજબ}\\
      & = \FV{P} \setminus \{y\} && \text{આગમનની ધારણાથી} \\
      & = \FV{\lambd[y][P]} && \text{\olref[fv]{def:fv}\olref[fv]{def:fv2} મુજબ}
    \end{aligned}
    \end{multline*}
  \end{enumerate}
\sourcecorrection{OLLAM-013}{સ્થિર અંગ્રેજી મૂળની આ સાબિતીમાં
$x\notin\FV{Q}$ લખાયું છે, પરંતુ આ કિસ્સામાં~$M$ એ
$\lambd[y][P]$ છે; જરૂરી ધારણા~$x\notin\FV{P}$ છે.
બાંધનાર ચલ જ બદલવાનું હોય તે યથાવત રહેતો કિસ્સો પણ
અલગ બતાવ્યો છે અને સમીકરણ-શ્રેણીમાંનું બેવડું સમાનચિહ્ન દૂર
કર્યું છે.}
\end{proof}

\begin{prob}
  \olref[lam][syn][sub]{thm:notinfv}ની સાબિતી પૂર્ણ કરો.
\end{prob}

\begin{thm} \ollabel{thm:infv}
  $x \in \FV{M}$ હોય અને સ્થાનાપન્ન વ્યાખ્યાયિત હોય તો
  $\FV{\Subst{M}{N}{x}} = (\FV{M} \setminus
  \{x\}) \cup \FV{N}$.
\sourcecorrection{OLLAM-014}{સ્થિર અંગ્રેજી મૂળની ધારણામાં
$x\in\FV{M})$ લખાયેલો વધારાનો બંધ કૌંસ દૂર કર્યો છે.}
\end{thm}

\begin{proof}
  $M$ના રચનાક્રમ પર આગમનથી.
  \begin{enumerate}
  \item $M$ ચલ હોય: અભ્યાસપ્રશ્ન.
  \item $M$ એ~$PQ$ સ્વરૂપનું હોય: $\Subst{(PQ)}{N}{x}$
    વ્યાખ્યાયિત હોવાથી તે
    $(\Subst{P}{N}{x})(\Subst{Q}{N}{x})$ છે, જ્યાં બંને
    સ્થાનાપન્ન વ્યાખ્યાયિત છે. વળી~$x \in \FV{PQ}$
    હોવાથી~$x \in \FV{P}$ અથવા~$x \in \FV{Q}$ અથવા
    બંને સાચાં છે. બાકીનો ભાગ અભ્યાસપ્રશ્ન છે.
    \sourcecorrection{OLLAM-015}{સ્થિર અંગ્રેજી મૂળમાં
    અનુપ્રયોગના આ કિસ્સામાં~$\Subst{(PQ)}{N}{y}$ છે,
    જ્યારે આખા પ્રમેયમાં બદલવાનો ચલ~$x$ છે. અહીં~$x$
    મૂક્યું છે.}
  \item $M$ એ~$\lambd[y][P]$ સ્વરૂપનું હોય.
    $\Subst{\lambd[y][P]}{N}{x}$ વ્યાખ્યાયિત હોવાથી તે
    $\lambd[y][\Subst{P}{N}{x}]$ છે; અહીં
    $\Subst{P}{N}{x}$ વ્યાખ્યાયિત છે, $x \neq y$ અને
    $y \notin \FV{N}$. વળી, $x \in
    \FV{\lambd[y][P]}$ હોવાથી~$x \in \FV{P}$ પણ છે.
    હવે:
    \sourcecorrection{OLLAM-016}{સ્થિર અંગ્રેજી મૂળમાં અહીં
    $y\in\FV{\lambd[x][P]}$ પરથી~$y\in\FV{P}$ કહે છે;
    આ કિસ્સાનું પદ તો~$\lambd[y][P]$ છે અને જરૂરી
    ધારણા~$x\in\FV{\lambd[y][P]}$, તેથી~$x\in\FV{P}$ છે.}
    \begin{multline*}
      \FV{\Subst{(\lambd[y][P])}{N}{x}} \\
      \begin{aligned}
      & = \FV{\lambd[y][\Subst{P}{N}{x}]} \\
      & = \FV{\Subst{P}{N}{x}} \setminus \{y\} \\
      & = ((\FV{P} \setminus \{x\}) \cup \FV{N}) \setminus \{y\}
       && \text{આગમનની ધારણાથી} \\
      & = (\FV{P} \setminus \{x, y\}) \cup \FV{N}
       && y \notin \FV{N} \\
      & = (\FV{\lambd[y][P]} \setminus \{x\}) \cup \FV{N}
      \end{aligned}
      \end{multline*}
    \sourcecorrection{OLLAM-017}{સ્થિર અંગ્રેજી મૂળની
    આગમન-શ્રેણીમાં~$\FV{P}\setminus\{y\}$ અને
    $\FV{N}\setminus\{x\}$ ખોટી જગ્યાએ છે તથા કૌંસ
    અધૂરા છે. અહીં પહેલા~$x$ દૂર કર્યા પછી સંઘમાંથી~$y$
    દૂર કર્યો છે; છેલ્લું સરળીકરણ~$y\notin\FV{N}$થી
    મળે છે, મૂળમાં લખેલા~$x\notin\FV{N}$થી નહીં.
    સમીકરણ-શ્રેણીનું બેવડું સમાનચિહ્ન પણ દૂર કર્યું છે.}
  \end{enumerate}
\end{proof}

\begin{prob}
  \olref[lam][syn][sub]{thm:infv}ની સાબિતી પૂર્ણ કરો.
\end{prob}

\begin{thm}\ollabel{thm:clr}
  સ્થાનાપન્ન વ્યાખ્યાયિત હોય અને~$x \notin \FV{N}$ હોય તો
  $x \notin \FV{\Subst{M}{N}{x}}$.
\end{thm}

\begin{proof}
  અભ્યાસપ્રશ્ન.
\end{proof}

\begin{prob}
  \olref[lam][syn][sub]{thm:clr} સાબિત કરો.
\end{prob}

\begin{thm}\ollabel{thm:inv}
  $\Subst{M}{y}{x}$ વ્યાખ્યાયિત હોય અને
  $y \notin \FV{M}$ હોય તો
  $\Subst{\Subst{M}{y}{x}}{x}{y} = M$.
\end{thm}

\begin{proof}
  $M$ના રચનાક્રમ પર આગમનથી.
  \begin{enumerate}
  \item $M$ ચલ~$z$ હોય: અભ્યાસપ્રશ્ન.
  \item $M$ એ~$(PQ)$ સ્વરૂપનું હોય. પછી:
    \begin{align*}
      \Subst{\Subst{(PQ)}{y}{x}}{x}{y}
      &=\Subst{((\Subst{P}{y}{x})(\Subst{Q}{y}{x}))}{x}{y} \\
      &= (\Subst{\Subst{P}{y}{x}}{x}{y})(\Subst{\Subst{Q}{y}{x}}{x}{y}) \\
      &= (PQ) \text{આગમનની ધારણાથી}
    \end{align*}
  \item $M$ એ~$\lambd[z][N]$ સ્વરૂપનું હોય. બદલવાના ચલની
    મુક્ત ઘટના જ ન હોય તો બંને સ્થાનાપન્ન પદને યથાવત રાખે છે.
    બાકીના કિસ્સામાં~$\Subst{\lambd[z][N]}{y}{x}$
    વ્યાખ્યાયિત હોવાથી~$z \neq y$ અને~$z\neq x$ છે.
    તેથી:
    \begin{align*}
      \Subst{\Subst{(\lambd[z][N])}{y}{x}}{x}{y}\\
      & = \Subst{(\lambd[z][\Subst{N}{y}{x}])}{x}{y} \\
      & = \lambd[z][\Subst{\Subst{N}{y}{x}}{x}{y}] \\
      &= \lambd[z][N] \text{આગમનની ધારણાથી}
    \end{align*}
    \sourcecorrection{OLLAM-018}{સ્થિર અંગ્રેજી મૂળ પહેલું
    સ્થાનાપન્ન વ્યાખ્યાયિત હોવાથી સીધું~$z\neq y$
    માને છે; બદલવાના ચલની મુક્ત ઘટના ન હોય તો આ જરૂરી
    નથી. તે યથાવત રહેતો કિસ્સો અલગ લીધો છે. બાકીનું
    પ્રદર્શન બંને બંધન-અથડામણો ટળે એવા~$z\neq x,y$
    કિસ્સામાં છે.}
  \end{enumerate}
\end{proof}

\begin{prob}
  \olref[lam][syn][sub]{thm:inv}ની સાબિતી પૂર્ણ કરો.
\end{prob}

\end{document}
